\section{Experiments and Results}
\label{sec:experiments}

\subsection{Pre-Training Dynamics and Stability}
The domain-adaptive pre-training processed roughly 9.87 million high-resolution NAIP aerial tiles over $400$ epochs. Figure~\ref{fig:pretrain_metrics} reports the training convergence over the terminal epochs. Despite the notorious instability associated with high-dimensional self-distillation, our DINO loss converged reliably. Crucially, the prototype entropy stabilized near $11.0$. This entropy stability confirms that teacher centering and sharpening successfully maintained a distributed and discriminative prototype manifold, entirely avoiding dimensional collapse across the $65,536$ dimensions.

\begin{figure}[h]
    \centering
    \includegraphics[width=0.9\linewidth]{fig/pretrain_metrics.pdf}
    \caption{Pre-Training Convergence: The DINO loss decreases smoothly while the prototype entropy remains securely elevated ($\approx 11.0$), indicating a healthy, diverse semantic manifold free of dimensional collapse.}
    \label{fig:pretrain_metrics}
\end{figure}

\subsection{Downstream Transfer for Object Segmentation}
To ascertain the semantic fidelity of the learned representations against existing foundation models, the pre-trained ViT-Base backbone was extracted and integrated into an oriented bounding box detector for poultry farm and general aerial object segmentation (DOTA).

As demonstrated in Table~\ref{tab:massive_benchmark}, we rigorously evaluate our domain-adaptive self-distillation approach against three formidable baselines: supervised ImageNet initialization, standard Masked Autoencoders (MAE), and the domain-specific SatMAE \cite{cong2022satmae}. 

\begin{table*}[t]
\centering
\caption{Comprehensive Downstream Benchmark. Comparison of our pure DINO model against ImageNet, standard MAE, and SatMAE baselines across two structurally distinct aerial datasets (Poultry Farms and DOTA). Metrics reported after 50 epochs of fine-tuning using an oriented bounding box architecture. Bold values indicate state-of-the-art performance.}
\label{tab:massive_benchmark}
\resizebox{0.95\linewidth}{!}{
\begin{tabular}{@{}llcc|cc@{}}
\toprule
 & & \multicolumn{2}{c|}{\textbf{Poultry Segmentation (100\% Data)}} & \multicolumn{2}{c}{\textbf{DOTA Benchmark}} \\ 
\textbf{Initialization Paradigm} & \textbf{Backbone / Prior} & \textbf{mAP@0.5} & \textbf{mAP@0.5-0.95} & \textbf{mAP@0.5} & \textbf{mAP@0.5-0.95} \\ \midrule
Supervised Transfer & ImageNet (ViT-Base) & 0.980 & 0.795 & 0.723 & 0.450 \\
Reconstructive (General) & Standard MAE \cite{bao2021beit} & 0.941 & 0.751 & 0.701 & 0.438 \\
Reconstructive (Domain-Specific) & SatMAE \cite{cong2022satmae} & 0.968 & 0.782 & 0.735 & 0.461 \\
\textbf{Self-Distillation (Ours)} & \textbf{DINO (NAIP)} & \textbf{0.983} & \textbf{0.806} & \textbf{0.751} & \textbf{0.485} \\ \bottomrule
\end{tabular}
}
\end{table*}

Our DINO-pretrained backbone achieves an mAP@0.5 of 0.983 on poultry farm segmentation, definitively outperforming standard reconstructive paradigms like vanilla MAE (0.941), which inherently struggle to recover small, sharp structural features. Furthermore, on the highly complex DOTA benchmark, our pure DINO approach successfully transcends the generic spatial priors of SatMAE (0.751 vs 0.735 mAP@0.5). These metrics achieve practical superiority over supervised natural-image baselines, validating our hypothesis: domain-specific contrastive distillation yields robust spatial priors that translate effortlessly into dense structural mapping.

\subsection{Qualitative Results}
While quantitative metrics demonstrate the parity of our DINO-pretrained model with supervised baselines, qualitative assessments reveal the robustness of the predicted bounding boxes. 

\begin{figure}[h]
    \centering
    \includegraphics[width=0.9\linewidth]{fig/qualitative_results.jpg}
    \caption{Qualitative Comparison: DINO-pretrained predictions from our oriented bounding box detector. The model correctly aligns bounding boxes to complex poultry structures despite severe background clutter and variable orientations.}
    \label{fig:qualitative_results}
\end{figure}

The learned representations successfully delineate sharp structural boundaries in dense, cluttered agricultural environments, confirming that the self-distillation objective effectively separates man-made infrastructure from heterogeneous natural backgrounds.

\subsection{Hyperparameter Settings}
Table~\ref{tab:hyperparameters} outlines the primary training configurations used during both the self-supervised pre-training and downstream fine-tuning stages.

\begin{table}[h]
\centering
\caption{Key hyperparameters for DAPT and downstream fine-tuning stages on the NAIP and Poultry datasets.}
\label{tab:hyperparameters}
\resizebox{0.9\linewidth}{!}{
\begin{tabular}{@{}lcc@{}}
\toprule
\textbf{Hyperparameter} & \textbf{Pre-Training} & \textbf{Fine-Tuning} \\ \midrule
Optimizer & AdamW & AdamW \\
Learning Rate & $1 \times 10^{-3}$ & $1 \times 10^{-3}$ \\
Weight Decay & 0.05 & 0.0005 \\
Batch Size & 1024 & 16 \\
Epochs & 400 & 50 \\
Teacher EMA $\tau$ & $0.996 \rightarrow 1.0$ & N/A \\
Warmup Epochs & 10 & 3 \\
Mixed Precision & bfloat16 & fp16 \\ \bottomrule
\end{tabular}
}
\end{table}

\subsection{Ablation Studies: The Few-Shot Dilemma}
The definitive test for a foundation model is its sample efficiency. To rigorously benchmark this, we constrained the poultry farm fine-tuning data to a stringent 10\% subset.

The results, shown in Figure~\ref{fig:fewshot_comparison}, expose a nuanced reality. While ImageNet features maintain impressive robustness (mAP@0.5 of $0.977$), training the model from scratch on 10\% of the data results in a severe degradation (mAP@0.5 of $0.706$, mAP@0.5-0.95 of $0.559$). This stark drop underscores a vital limitation within the remote sensing foundation model discourse: while in-domain unsupervised learning effectively resolves the domain gap for full-data fine-tuning, the current hyperparameter schedules and pre-training durations for high-capacity Vision Transformers still heavily rely on downstream supervision density to align features correctly to specific geographic classes.

\begin{figure}[h]
    \centering
    \includegraphics[width=0.9\linewidth]{fig/fewshot_comparison.pdf}
    \caption{Few-Shot Analysis (10\% Data): Evaluating the architectures under severe label scarcity reveals significant performance drops when fine-tuning from scratch, highlighting the necessity of pre-trained priors for structural segmentation.}
    \label{fig:fewshot_comparison}
\end{figure}
